\par\bigskip\pdfbookmark[0]{De nexu. Fragmentum V}{nexus-fragment-V}\centerline{[V.]}\medskip
Multitudo classium mediocris\begingroup\renewcommand{\thefootnote}{*)}\footnote{[Vergl. \emph{Disquiss. Arithm.} art. 302; die dortige Formel weicht um eine Constante $\delta$ von der hier im Text vorkommenden ab.]}\endgroup\ circa determinantem negativum $-D$ est proxime
\[
=\frac{\pi\surd D}{4(1+\tfrac1{27}+\tfrac1{125}+\text{etc.})}
\]
Multitudo vera exprimitur sequentibus formulis, ubi brevitatis caussa scribitur $m$ pro multitudine mediocri, $M$ pro vera; $p,q$ exprimunt omnes numeros impares primos ipsum $D$ non metientes, ille divisores, hic non-divisores ipsius $\square+D$; $r$ numeros\begingroup\renewcommand{\thefootnote}{**)}\footnote{[impares.]}\endgroup\ primos ipsum $D$ metientes:

\SourceLeaf{295}
\[
\begin{aligned}
\text{I.}\quad M&=m\ \text{Prod. ex }\frac{p^3+p^2}{p^3-1}\pdot\frac{q^3-q^2}{q^3-1}\pdot\frac{r^3-r}{r^3-1}\\
\text{II.}\quad M&=\frac{\pi\surd D}4\ \text{Prod. ex }\frac{p+1}p\pdot\frac{q-1}q\pdot\frac{rr-1}{rr}\\
\text{III. NB.}\quad M&=\frac{2\surd D}\pi\ \text{Prod. ex }\frac p{p-1}\pdot\frac q{q+1}\\
\text{IV.}\quad M&=\surd\left\{\frac D2\pdot\text{Prod. ex }\frac{p+1}{p-1}\pdot\frac{q-1}{q+1}\pdot\frac{rr-1}{rr}\right\}\\
\text{V.}\quad M&=\frac{2\surd D}\pi\{1\pm\tfrac13\pm\tfrac15\pm\tfrac17\pm\tfrac19\text{ etc.}\}
\end{aligned}
\]
NB. Die Formel III wird unmittelbar aus der Vergleichung der beiden Arten, die darstellbaren Zahlen bis zu einer gewissen Grenze zu zählen, abgeleitet.
\BookRule{.30\textwidth}
